% Hop-area is collinear with CHM energy. Not a second Jacobson term.
\documentclass[11pt]{article}
\usepackage[margin=1.05in]{geometry}
\usepackage{amsmath,amssymb,amsthm}
\usepackage[colorlinks=true,linkcolor=blue,citecolor=blue,urlcolor=blue]{hyperref}

\newtheorem*{admission}{Admission}

\title{\textbf{Hop-Area Is an Energy Proxy,\\
Not a Second First-Law Term}}
\author{Robert W.\ McGwier\\
\small CTO, Cohere Technology Group\\
\small GitHub: \texttt{n4hy}}
\date{15 August 2026 --- Version 1}

\begin{document}
\maketitle

\begin{abstract}
\noindent
Jacobson and FGHMV need two pieces: $\delta S=\eta\,\delta A
+\beta\,\delta E$. A conformal hop bump changes both.
This note splits the variations.

Matter only, hops fixed: $\delta A\equiv 0$,
$\rho(\delta S,P_{\mathrm{CHM}})=0.9998$. The first law
does not need area.

Geometry only: $\rho(\delta A,P)=0.996$. Area and energy
are the same bump. A raw partial correlation of $-0.987$
is the degeneracy artifact of that collinearity and is
not scored as independence.

Auditor, $N=10$: $\rho(A,P)=0.988$, independence REFUTED.

Papers 39--41 saw $\delta S$ track area because they saw
it track energy. Not $8\pi G$. Not Einstein. Not de~Sitter.
\end{abstract}

\begin{admission}
I asked whether area was a second term. It is not. I did
not promote a two-column residual drop under
$\rho=0.996$ into Jacobson.
\end{admission}

\end{document}
